\documentclass[11pt]{article}

\usepackage[margin=1in]{geometry}
\usepackage{amsmath,amssymb,amsthm,mathtools}
\usepackage{booktabs,longtable,array}
\usepackage{microtype}
\usepackage[hidelinks]{hyperref}

\newtheorem{theorem}{Theorem}[section]
\newtheorem{lemma}[theorem]{Lemma}
\newtheorem{proposition}[theorem]{Proposition}
\newtheorem{corollary}[theorem]{Corollary}
\theoremstyle{remark}
\newtheorem{remark}[theorem]{Remark}

\newcommand{\C}{\mathbb C}
\newcommand{\D}{\mathbb D}
\newcommand{\e}{\mathrm e}
\newcommand{\Rea}{\operatorname{Re}}

\title{Sendov's Conjecture in Degrees Nine, Ten, and Eleven}
\author{Anton Shakov}
\date{}

\begin{document}
\maketitle

\begin{abstract}
We prove Sendov's conjecture for complex polynomials of degrees nine, ten, and
 eleven.  The three proofs use the same reciprocal normalization.  A
 hypothetical counterexample yields a product constraint, a localization
 inequality for the mean reciprocal critical-point displacement, and an exact
 integral identity.  After centering the reciprocal critical points, the
 analytic problem reduces to positivity of explicit bivariate polynomials.
 Degree nine requires eighteen interior Bernstein certificates and one
 boundary certificate; degree ten requires fifty interior certificates and one
 boundary certificate; degree eleven requires a split-integral estimate and
 530 interior certificates together with one boundary certificate.  Every
 finite calculation is carried out in exact rational arithmetic, with an
 independent symbolic reconstruction for each degree.
\end{abstract}

\section{Introduction}

Sendov's conjecture asserts that if all zeros of a complex polynomial lie in
 the closed unit disk, then the closed unit disk centered at each zero contains
 a zero of the derivative.  Brown and Xiang proved the conjecture for degrees
 at most eight \cite{BrownXiang}.  We establish the next three degrees.

\begin{theorem}\label{thm:main}
Let
\[
 p(z)=\prod_{k=1}^{n}(z-z_k),\qquad |z_k|\le 1,
\]
where $n\in\{9,10,11\}$.  All zeros are counted with multiplicity.  For every zero $a$ of $p$, there is a zero $\zeta$
 of $p'$ such that
\[
 |a-\zeta|\le 1.
\]
\end{theorem}

By rotation, the distinguished zero may be taken to be a real number
$a\in[0,1]$.  We argue by contradiction and call the polynomial \emph{bad at
$a$} if every zero of $p'$ has distance strictly greater than $1$ from $a$.
The proof is computer-assisted only in finite polynomial-positivity checks.
The reductions, the certificate polynomials, and the exact verification
procedure are all given below.

\section{Preliminaries}

\subsection{Separation of the remaining zeros}

We use the Grace--Walsh--Szeg\H{o} coincidence theorem in its standard form for
symmetric multi-affine polynomials on disks; see, for example,
\cite{BrandenWagner}.

\begin{lemma}\label{lem:separation}
Let $P$ be a polynomial of degree $n\ge 2$, let $\alpha$ be a zero of $P$, and
suppose that
\[
 |\alpha-\zeta|>\rho
\]
for every zero $\zeta$ of $P'$.  Then $\alpha$ is simple and every other zero
$\beta$ of $P$ satisfies
\[
 |\beta-\alpha|>2\rho\sin\frac{\pi}{n}.
\]
\end{lemma}

\begin{proof}
The simplicity of $\alpha$ is immediate.  Let $\zeta_1,\ldots,\zeta_{n-1}$ be
 the zeros of $P'$ and put $w=\beta-\alpha$.  Since
$P(\beta)=P(\alpha)$,
\[
 0=\int_0^1\frac{P'(\alpha+tw)}{P'(\alpha)}\,dt
   =\int_0^1\prod_{j=1}^{n-1}(1-tu_j)\,dt,
 \qquad
 u_j=-\frac{w}{\alpha-\zeta_j}.
\]
The symmetric multi-affine polynomial
\[
 F(u_1,\ldots,u_{n-1})=\int_0^1\prod_{j=1}^{n-1}(1-tu_j)\,dt
\]
has diagonal
\[
 F(u,\ldots,u)=\frac{1-(1-u)^n}{nu},
\]
with the value $1$ at $u=0$.  Its diagonal zeros are
$1-\e^{2\pi i k/n}$, $1\le k<n$, and therefore have modulus at least
$2\sin(\pi/n)$.  The coincidence theorem implies that $F$ is nonzero when
all $|u_j|<2\sin(\pi/n)$.  If
$|w|\le 2\rho\sin(\pi/n)$, then the strict critical-point separation gives
this inequality for every $u_j$, a contradiction.
\end{proof}

\subsection{A reciprocal-square extremum}

\begin{lemma}\label{lem:sigma}
Let $0<m<M$, let $r_1,\ldots,r_N\in[m,M]$, and suppose
\[
 \prod_{k=1}^N r_k\ge C>m^N.
\]
Let $\nu$ be the least integer for which
$M^\nu m^{N-\nu}\ge C$.  Then
\[
 \sum_{k=1}^N\frac1{r_k^2}
 \le
 \frac{N-\nu}{m^2}+\frac{\nu-1}{M^2}
 +\left(\frac{M^{\nu-1}m^{N-\nu}}{C}\right)^2.
\]
\end{lemma}

\begin{proof}
Put $r_k=\e^{t_k}$.  Since $t\mapsto\e^{-2t}$ is decreasing, we may reduce
$\sum t_k$ to $\log C$.  On the resulting section of the box
$[\log m,\log M]^N$, the convex function $\sum\e^{-2t_k}$ is maximized at an
extreme point.  Hence all but at most one coordinate are endpoints.  By the
choice of $\nu$, the maximizing configuration has $N-\nu$ entries equal to
$m$, $\nu-1$ entries equal to $M$, and one entry equal to
$C/(M^{\nu-1}m^{N-\nu})$.
\end{proof}

\subsection{Centered elementary symmetric functions}

For $d=(d_1,\ldots,d_N)$, let $e_m(d)$ denote the $m$th elementary symmetric
function.

\begin{lemma}[variance under differentiation]\label{lem:variance-derivative}
Let $f$ be a monic polynomial of degree $N\ge2$ whose zeros
$d_1,\ldots,d_N$ have mean zero, and let $w_1,\ldots,w_{N-1}$ be the zeros of
$f'$.  Then
\[
 \sum_{j=1}^{N-1}|w_j|^2
 \le \frac{N-2}{N}\sum_{j=1}^N|d_j|^2.
\]
\end{lemma}

\begin{proof}
Let $A=\operatorname{diag}(d_1,\ldots,d_N)$,
$u=N^{-1/2}(1,\ldots,1)^{\mathsf T}$, and $P=I-uu^*$.  The compression
$B=PAP|_{u^\perp}$ has characteristic polynomial $f'/N$.  Indeed, the Schur
complement identity gives, first away from the zeros of $f$ and then
identically,
\[
 \det(zI_{u^\perp}-B)
 =f(z)\,u^*(zI-A)^{-1}u
 =\frac{f'(z)}N.
\]
By Schur triangularization,
\[
 \sum_{j=1}^{N-1}|w_j|^2\le \|B\|_{\mathrm{HS}}^2.
\]
Since $u^*Au=0$ and $A$ is normal,
\[
 \|B\|_{\mathrm{HS}}^2
 =\|PAP\|_{\mathrm{HS}}^2
 =\|A\|_{\mathrm{HS}}^2-\|Au\|^2-\|A^*u\|^2
 =\frac{N-2}{N}\sum_{j=1}^N|d_j|^2.
\]
\end{proof}

\begin{lemma}\label{lem:centered-esp}
Suppose
\[
 \sum_{j=1}^N d_j=0,
 \qquad
 \frac1N\sum_{j=1}^N|d_j|^2\le \eta^2.
\]
Then
\begin{align}
 |e_2(d)|&\le \frac N2\eta^2,\label{eq:e2}\\
 |e_3(d)|&\le \frac{N\sqrt{N-1}}3\eta^3,\label{eq:e3}\\
 |e_m(d)|&\le
 \left(\frac{N^N}{m^m(N-m)^{N-m}}\right)^{1/2}\eta^m
 \quad(1\le m<N),\label{eq:cauchy-esp}\\
 |e_m(d)|&\le \binom Nm
 \left(\frac{m-1}{N-1}\right)^{m/2}\eta^m
 \quad(2\le m\le N),\label{eq:schoenberg}\\
 |e_N(d)|&\le \eta^N.\label{eq:topesp}
\end{align}
\end{lemma}

\begin{proof}
Newton's identities give $e_2=-\frac12\sum d_j^2$, proving
\eqref{eq:e2}.  If $V=\sum|d_j|^2$, then
\[
 \max_j|d_j|^2\le\frac{N-1}{N}V,
\]
because $d_j=-\sum_{k\ne j}d_k$.  Since
$e_3=\frac13\sum d_j^3$,
\[
 |e_3|\le\frac13\max_j|d_j|\,V
 \le\frac13\sqrt{\frac{N-1}{N}}V^{3/2},
\]
which is \eqref{eq:e3}.

For \eqref{eq:cauchy-esp}, let
$H(z)=\prod_{j=1}^N(1+zd_j)$.  On $|z|=r$, the arithmetic--geometric mean
inequality and $\sum d_j=0$ give
\[
 |H(z)|
 \le\left(\frac1N\sum_{j=1}^N|1+zd_j|^2\right)^{N/2}
 \le(1+r^2\eta^2)^{N/2}.
\]
Cauchy's coefficient estimate yields
\[
 |e_m(d)|\le r^{-m}(1+r^2\eta^2)^{N/2}.
\]
For $1\le m<N$, minimization at
$r^2=m/((N-m)\eta^2)$ proves \eqref{eq:cauchy-esp}; the case $\eta=0$ is
immediate.

For \eqref{eq:schoenberg}, let $w_1,\ldots,w_m$ be the zeros of
$f^{(N-m)}$, where $f(z)=\prod_{j=1}^N(z-d_j)$.  Repeated application of
Lemma~\ref{lem:variance-derivative} gives
\[
 \frac1m\sum_{j=1}^m|w_j|^2
 \le\frac{m-1}{N-1}\eta^2.
\]
Coefficient comparison gives
\[
 \prod_{j=1}^m w_j=\frac{e_m(d)}{\binom Nm}.
\]
The arithmetic--geometric mean inequality proves \eqref{eq:schoenberg}.
Finally, \eqref{eq:topesp} is the arithmetic--geometric mean inequality
applied directly to $|d_1|^2,\ldots,|d_N|^2$.
\end{proof}

\subsection{Bernstein positivity}

\begin{lemma}\label{lem:bernstein}
Let $R(U,V)$ have bidegree at most $(r,s)$ and write
\[
 R(U,V)=\sum_{i=0}^r\sum_{j=0}^s b_{ij}
 \binom riU^i(1-U)^{r-i}
 \binom sjV^j(1-V)^{s-j}.
\]
If every $b_{ij}>0$, then $R>0$ on $[0,1]^2$.
\end{lemma}

\begin{proof}
The tensor-product Bernstein basis functions are nonnegative and sum to $1$.
\end{proof}

\section{The reciprocal-integral reduction}

Let $n\ge2$ and put $N=n-1$.  Write
\[
 p(z)=(z-a)\prod_{k=1}^N(z-z_k),
 \qquad
 p'(z)=n\prod_{j=1}^N(z-\zeta_j).
\]
Assume that $p$ is bad at $a$.  Then $a$ is simple.  Set
\[
 X_k=\frac1{a-z_k},\qquad
 Y_j=\frac1{a-\zeta_j},\qquad
 \mu=\frac1N\sum_{j=1}^NY_j=u+iv.
\]
Evaluation of $p'$ and $p''/p'$ at $a$ gives
\begin{equation}\label{eq:anchors-general}
 \prod_{k=1}^N(a-z_k)=n\prod_{j=1}^N(a-\zeta_j),
 \qquad
 \sum_{j=1}^NY_j=2\sum_{k=1}^NX_k.
\end{equation}
Thus
\begin{equation}\label{eq:mu-general}
 \sum_{k=1}^NX_k=\frac N2\mu,
 \qquad
 |\mu|<1.
\end{equation}

There is also a full coefficient identity.  If
\[
 F(t)=\prod_{k=1}^N(1-tX_k),
 \qquad
 G(t)=\prod_{j=1}^N(1-tY_j),
\]
then
\begin{equation}\label{eq:F-G}
 G(t)=F(t)+tF'(t),
 \qquad
 e_m(Y)=(m+1)e_m(X)\quad(0\le m\le N).
\end{equation}
Indeed, $F(t)=q(a-t)/q(a)$ for $p(z)=(z-a)q(z)$, while
$G(t)=p'(a-t)/p'(a)$.

Put
\[
 r_k=|a-z_k|,
 \qquad
 \sigma=\sum_{k=1}^N r_k^{-2}.
\]
Lemma~\ref{lem:separation}, the unit-disk hypothesis, and
\eqref{eq:anchors-general} give
\begin{equation}\label{eq:annulus-product-general}
 2\sin\frac\pi n<r_k\le1+a,
 \qquad
 \prod_{k=1}^Nr_k>n.
\end{equation}
Moreover, since $z_k=a-X_k^{-1}$,
\[
 2a\Rea X_k+(1-a^2)|X_k|^2
 =1+(1-|z_k|^2)|X_k|^2\ge1.
\]
Summing and using \eqref{eq:mu-general}, we obtain
\begin{equation}\label{eq:localization-general}
 aNu+(1-a^2)\sigma\ge N.
\end{equation}

Set
\[
 \eta^2=1-|\mu|^2,
 \qquad
 D_j=Y_j-\mu.
\]
Then
\begin{equation}\label{eq:variance-general}
 \sum_{j=1}^ND_j=0,
 \qquad
 \sum_{j=1}^N|D_j|^2\le N\eta^2.
\end{equation}
Choose constants $c_m$ such that
\begin{equation}\label{eq:cm-general}
 |e_m(D)|\le c_m\eta^m
 \qquad(2\le m\le N).
\end{equation}

Finally, define
\begin{equation}\label{eq:I-general}
 I=\int_0^aG(t)\,dt=-\frac{p(0)}{p'(a)}.
\end{equation}
Since $|p(0)|\le a$ and $|p'(a)|>n$,
\begin{equation}\label{eq:I-upper-general}
 |I|<\frac an.
\end{equation}

We now state the polynomial criterion used in all three degrees.  Suppose that
on an interval $[\alpha,\beta]$ we have $\sigma\le S$.  Put
\begin{align}
 L_S(a)&=\frac{N-S(1-a^2)}{aN},\label{eq:LS}\\
 q_0(a,\eta;v)&=
 1-\frac{2(N-S+Sa^2)}N v+a^2(1-\eta^2)v^2,\label{eq:q0}\\
 h(a,\eta)&=
 \left(\frac{2S}{N}-1\right)(1-a^2)-a^2\eta^2,
 \label{eq:h-general}\\
 \ell_m&=\left\lfloor\frac{N-m}{2}\right\rfloor,
 \qquad
 k=\left\lfloor\frac n2\right\rfloor.
\end{align}
For rational $S,\lambda$, define
\begin{equation}\label{eq:Pfull}
\begin{split}
 P_{n,S,\lambda}(a,\eta)
 ={}&\frac{1-a}{n}+\frac{\eta^2}{2n}
 -\frac{h(a,\eta)^k}{n\lambda}\\
 &-\sum_{m=2}^N c_m a^{m+1}\eta^m
 \int_0^1v^m q_0(a,\eta;v)^{\ell_m}\,dv.
\end{split}
\end{equation}
This is a polynomial with rational coefficients.

\begin{proposition}[unsplit certificate]\label{prop:full-certificate}
Assume that, for $a\in[\alpha,\beta]$,
\begin{equation}\label{eq:box-conditions}
 \sigma\le S,
 \qquad
 L_S(a)\ge\lambda>0,
 \qquad
 Y^2\ge1-\lambda^2,
 \qquad
 2\lambda\ge\beta,
\end{equation}
and
\begin{equation}\label{eq:h-condition}
 0\le\frac{2S}{N}-1,
 \qquad
 \left(\frac{2S}{N}-1\right)(1-\alpha^2)\le1.
\end{equation}
If
\[
 P_{n,S,\lambda}(a,\eta)>0
 \qquad
 (\alpha\le a\le\beta,\ 0\le\eta\le Y),
\]
then no bad polynomial exists with distinguished zero in
$[\alpha,\beta]$.
\end{proposition}

\begin{proof}
By \eqref{eq:localization-general}, $u\ge L_S(a)\ge\lambda$, and hence
$0\le\eta^2\le1-\lambda^2\le Y^2$.  Put
\[
 J=\int_0^a(1-t\mu)^N\,dt
   =\frac{1-(1-a\mu)^n}{n\mu}.
\]
For $t=av$,
\[
 |1-t\mu|^2
 \le1-2tL_S(a)+(1-\eta^2)t^2=q_0(a,\eta;v).
\]
The right side is nonnegative, and \eqref{eq:box-conditions} gives
$q_0\le1$.  Expanding about $\mu$,
\[
 G(t)=(1-t\mu)^N+
 \sum_{m=2}^N(-t)^me_m(D)(1-t\mu)^{N-m}.
\]
Since $0\le q_0\le1$, rounding every half-integral exponent down gives
\begin{equation}\label{eq:error-full}
 |I-J|\le
 \sum_{m=2}^N c_m a^{m+1}\eta^m
 \int_0^1v^m q_0^{\ell_m}\,dv.
\end{equation}

Let $R=|1-a\mu|$.  Then $R^2\le h(a,\eta)$.  On a feasible pair,
$h\ge0$, while \eqref{eq:h-condition} gives $h\le1$.  Thus $R\le1$ and
$R^n\le h^k$.  Since $|\mu|\ge\lambda$,
\[
 |J|\ge\frac{1-R^n}{n|\mu|}
 \ge\frac1n+\frac{\eta^2}{2n}-\frac{h^k}{n\lambda}.
\]
Together with \eqref{eq:error-full}, this yields
\[
 |I|-\frac an\ge P_{n,S,\lambda}(a,\eta)>0,
\]
contrary to \eqref{eq:I-upper-general}.
\end{proof}

For degree eleven we use a split version.  Here $N=10$.  Define
\begin{equation}\label{eq:q1}
 q_1(a;v)=1-\frac{2(N-S+Sa^2)}N v+a^2v^2.
\end{equation}
For $0\le\tau\le1$, set
\begin{equation}\label{eq:Psplit}
\begin{split}
 P_{n,S,\lambda}^{(\tau)}(a,\eta)
 ={}&\frac{1-a}{n}+\frac{\eta^2}{2n}
 -\frac{h(a,\eta)^k}{n\lambda}\\
 &-\sum_{m=2}^N c_m a^{m+1}\eta^m
 \int_0^\tau v^m q_0(a,\eta;v)^{\ell_m}\,dv\\
 &-a\int_\tau^1\bigl(q_0(a,\eta;v)^5+q_1(a;v)^5\bigr)\,dv.
\end{split}
\end{equation}

\begin{proposition}[split certificate]\label{prop:split-certificate}
Let $n=11$.  Under \eqref{eq:box-conditions} and
\eqref{eq:h-condition}, if
\[
 P_{11,S,\lambda}^{(\tau)}(a,\eta)>0
\]
throughout a rectangle contained in
$[\alpha,\beta]\times[0,Y]$, then that rectangle contains no feasible pair
arising from a bad polynomial.
\end{proposition}

\begin{proof}
On $0\le v\le\tau$, use the termwise estimate from
Proposition~\ref{prop:full-certificate}.  On $\tau\le v\le1$,
\[
 |G(av)|
 \le\left(\frac1N\sum_{j=1}^N|1-avY_j|^2\right)^{N/2}
 \le q_1(a;v)^5
\]
by the arithmetic--geometric mean inequality, while
$|(1-av\mu)^N|\le q_0(a,\eta;v)^5$.  Thus the last line of
\eqref{eq:Psplit} bounds the remaining part of $|I-J|$.  The estimate for
$J$ is unchanged.
\end{proof}

\section{Degree nine}

Here $N=8$.  We use
\[
 m_9=\frac{681}{1000}<2\sin\frac\pi9.
\]
Indeed, if $T=2\sin(\pi/9)$, then $3T-T^3=\sqrt3$, the function
$3t-t^3$ is increasing on $[0,1]$, and
\[
 3-(3m_9-m_9^3)^2
 =\frac{16853534459219919}{10^{18}}>0.
\]
The coefficient majorants are
\begin{equation}\label{eq:c9}
 (c_2,\ldots,c_8)=
 \left(4,\frac{264}{35},24,56,28,8,1\right).
\end{equation}
The first three follow from Newton's identities and
Lemma~\ref{lem:centered-esp}; for $m\ge5$, Maclaurin's inequality applied to
$|D_1|,\ldots,|D_8|$ gives $|e_m(D)|\le\binom8m\eta^m$.

\subsection{The interval $0\le a\le9/20$}

At $a=0$, the product inequality in
\eqref{eq:annulus-product-general} is impossible.  For
$0<a\le9/20$, apply Lemma~\ref{lem:sigma} with
$M=29/20$, $m=m_9$, and $C=9$.  The resulting bound is
\[
 \sigma<\frac{1101}{200}=:S_0.
\]
By \eqref{eq:localization-general},
\[
 u>\frac{8-S_0(1-a^2)}{8a}.
\]
The right side exceeds $1$ on this interval because
\[
 S_0a^2-8a+8-S_0>0;
\]
the quadratic is decreasing there and has value $781/80000$ at $a=9/20$.
This contradicts $|\mu|<1$.

\subsection{The interval $9/20\le a\le9/10$}

For each row of Table~\ref{tab:nine}, Lemma~\ref{lem:sigma}, with
$M=1+\beta$, gives $\sigma\le S$.  The remaining hypotheses of
Proposition~\ref{prop:full-certificate} are exact rational inequalities.

\begin{longtable}{@{}c c c c@{}}
\caption{Degree-nine interior certificate data.}\label{tab:nine}\\
\toprule
$[\alpha,\beta]$ & $S$ & $\lambda$ & $Y$\\
\midrule
\endfirsthead
\toprule
$[\alpha,\beta]$ & $S$ & $\lambda$ & $Y$\\
\midrule
\endhead
$[9/20,19/40]$ & $11043/2000$ & $49/50$ & $199/1000$\\
$[19/40,1/2]$ & $2783/500$ & $239/250$ & $587/2000$\\
$[1/2,21/40]$ & $2257/400$ & $1863/2000$ & $91/250$\\
$[21/40,11/20]$ & $1151/200$ & $1811/2000$ & $849/2000$\\
$[11/20,23/40]$ & $11819/2000$ & $879/1000$ & $477/1000$\\
$[23/40,3/5]$ & $764/125$ & $1699/2000$ & $66/125$\\
$[3/5,5/8]$ & $637/100$ & $817/1000$ & $577/1000$\\
$[5/8,13/20]$ & $13093/2000$ & $401/500$ & $239/400$\\
$[13/20,27/40]$ & $13113/2000$ & $81/100$ & $1173/2000$\\
$[27/40,7/10]$ & $3289/500$ & $409/500$ & $1151/2000$\\
$[7/10,29/40]$ & $529/80$ & $413/500$ & $141/250$\\
$[29/40,3/4]$ & $6661/1000$ & $1669/2000$ & $1103/2000$\\
$[3/4,31/40]$ & $1681/250$ & $843/1000$ & $269/500$\\
$[31/40,4/5]$ & $1701/250$ & $213/250$ & $131/250$\\
$[4/5,33/40]$ & $2761/400$ & $1723/2000$ & $127/250$\\
$[33/40,17/20]$ & $14041/2000$ & $109/125$ & $49/100$\\
$[17/20,7/8]$ & $7161/1000$ & $221/250$ & $187/400$\\
$[7/8,9/10]$ & $3663/500$ & $359/400$ & $883/2000$\\
\bottomrule
\end{longtable}

After the affine substitution
\[
 a=\alpha+(\beta-\alpha)U,
 \qquad
 \eta=YV,
\]
the eighteen polynomials \eqref{eq:Pfull} have bidegree at most $(9,8)$.
All $1620$ exact Bernstein coefficients are positive.  The least coefficient
is greater than $1/2100$.  Proposition~\ref{prop:full-certificate} excludes
the entire interval.

\subsection{The interval $9/10\le a\le1$}

At $a=1$, \eqref{eq:localization-general} gives $u\ge1$.  Assume
$9/10\le a<1$ and put $\delta=1-a$.  Lemma~\ref{lem:sigma}, with $M=2$,
gives
\[
 \sigma<\frac{39}{5}.
\]
Consequently
\[
 1-u\le\frac{19}{20}\delta,
 \qquad
 \eta^2\le\frac{19}{10}\delta-\frac{361}{400}\delta^2.
\]
Write
\[
 \delta=x^2,
 \qquad
 \eta=xy.
\]
Then
\[
 0<x<\frac8{25},
 \qquad
 0\le y<\frac75.
\]
For $0\le t\le1-x^2$, put
\[
 q(x,y;t)=(1-t)^2+x^2t\left(\frac{19}{10}-y^2t\right).
\]
The preceding estimate gives $|1-t\mu|^2\le q\le1$.  Define
\[
 \mathcal E(x,y)=
 \sum_{m=2}^8c_mx^{m-2}y^m
 \int_0^{1-x^2}t^m q(x,y;t)^{\lfloor(8-m)/2\rfloor}\,dt.
\]
The centered expansion gives $|I-J|\le\delta\mathcal E$, where
$J=\int_0^a(1-t\mu)^8dt$.  Moreover,
\[
 |J|\ge\frac19+\frac{\eta^2}{18}-\frac\delta{1000}.
\]
The last estimate follows from
$|1-a\mu|^2\le(19/10)\delta$, $|\mu|\ge181/200$, and the exact integer
inequality
\[
 200^2 19^9<1629^2 10^{10}.
\]
Thus
\[
 |I|-\frac a9\ge\delta\mathcal P(x,y),
 \qquad
 \mathcal P(x,y)=\frac19+\frac{y^2}{18}-\frac1{1000}-\mathcal E(x,y).
\]
The polynomial $\mathcal P$ has bidegree $(24,8)$.  After
$x=(8/25)U$, $y=(7/5)V$, all $225$ exact Bernstein coefficients exceed
$1/100$.  This contradicts \eqref{eq:I-upper-general} and completes degree
nine.

\section{Degree ten}

Here $N=9$.  We use
\[
 m_{10}=\frac{309}{500}<2\sin\frac\pi{10}
 =\frac{\sqrt5-1}{2},
\]
which is equivalent to $(2m_{10}+1)^2<5$.  The coefficient majorants are
\begin{equation}\label{eq:c10}
\begin{array}{c|cccccccc}
 m&2&3&4&5&6&7&8&9\\ \hline
 c_m&9/2&9&2201/100&2201/100&877/50&217/20&481/100&1.
\end{array}
\end{equation}
The cases $m=2,3,9$ follow from Lemma~\ref{lem:centered-esp}; the remaining
rational numbers dominate the Cauchy bound \eqref{eq:cauchy-esp}, as checked by exact
squaring.

For $0\le a\le29/100$,
\[
 \prod_{k=1}^9r_k\le\left(\frac{129}{100}\right)^9<10,
\]
contrary to \eqref{eq:annulus-product-general}.  On
$[29/100,2/5]$, Lemma~\ref{lem:sigma} with $M=7/5$ has $\nu=9$ and gives
\[
 S_0=8\left(\frac57\right)^2
 +\left(\frac{(7/5)^8}{10}\right)^2.
\]
The quadratic
\[
 S_0a^2-9a+9-S_0
\]
has positive Bernstein coefficients on $[29/100,2/5]$.  Hence
\eqref{eq:localization-general} gives $u>1$.

It remains to cover $[2/5,9/10]$.  For $i=40,\ldots,89$, set
\[
 \alpha_i=\frac{i}{100},
 \qquad
 \beta_i=\frac{i+1}{100},
\]
and let $S_i$ be the bound from Lemma~\ref{lem:sigma} with
$M=1+\beta_i$, $m=m_{10}$, and $C=10$.  Define $\lambda_i$ to be the largest
multiple of $10^{-5}$ for which
\[
 S_i a^2-9\lambda_i a+9-S_i\ge0
 \qquad(\alpha_i\le a\le\beta_i),
\]
certified by its degree-two Bernstein coefficients.  Let $Y_i$ be the least
multiple of $10^{-5}$ satisfying
$Y_i^2\ge1-\lambda_i^2$.

For each $i$, all hypotheses of
Proposition~\ref{prop:full-certificate} hold, and the exact Bernstein
coefficients of
\[
 P_{10,S_i,\lambda_i}(a,\eta)
\]
are positive on
$[\alpha_i,\beta_i]\times[0,Y_i]$.  These polynomials have bidegree at most
$(10,10)$.  Across all fifty rectangles, the least coefficient is greater
than $1/3100$.

For the boundary interval, take the bound $S_b$ from
Lemma~\ref{lem:sigma} with $M=2$ and set
\[
 \lambda_b=\frac{22}{25}.
\]
The hypotheses of Proposition~\ref{prop:full-certificate} hold on
$[9/10,1]$.  Moreover, localization gives
\[
 1-u\le C_b(1-a),
 \qquad
 C_b=\frac{S_b(1+9/10)-9}{9(9/10)},
\]
and hence $\eta^2\le2C_b(1-a)$.  Put $1-a=x^2$ and $\eta=xY$.  The exact
inequalities
\[
 \left(\frac8{25}\right)^2\ge\frac1{10},
 \qquad
 \left(\frac{31}{20}\right)^2\ge2C_b
\]
show that every feasible pair lies in
\[
 0\le x\le\frac8{25},
 \qquad
 0\le Y\le\frac{31}{20}.
\]
Divide $P_{10,S_b,\lambda_b}(1-x^2,xY)$ by $x^2$.
The quotient is a rational polynomial, and all its exact Bernstein
coefficients are positive; the least exceeds $7/100$.  The case $a=1$ again
follows directly from \eqref{eq:localization-general}.  This proves degree
ten.

\section{Degree eleven}

Here $N=10$.  We use
\[
 m_{11}=\frac{563}{1000}<2\sin\frac\pi{11}.
\]
Put $x=571/2000$.  Machin's identity and the alternating series for
$\arctan$ give
\[
 \pi=16\arctan\frac15-4\arctan\frac1{239}
 >16\left(\frac15-\frac1{3\cdot5^3}\right)-\frac4{239}
 =\frac{281476}{89625}
 >\frac{6281}{2000}=11x.
\]
Thus $x<\pi/11$.  Since $\sin t>t-t^3/6$ for $0<t<1$,
\[
 2\sin\frac\pi{11}
 >2\left(x-\frac{x^3}{6}\right)
 =\frac{563}{1000}+\frac{5830589}{24000000000}
 >\frac{563}{1000}.
\]  The coefficient majorants are
\begin{equation}\label{eq:c11}
\begin{array}{c|ccccccccc}
 m&2&3&4&5&6&7&8&9&10\\ \hline
 c_m&5&10&70/3&32&3125/108&2121/100&3125/256&509/100&1.
\end{array}
\end{equation}
Here $c_4$ is the bound \eqref{eq:schoenberg}; $c_5,c_6,c_8$ are the exact
Cauchy bounds; $c_7,c_9$ are rational upper bounds for
\eqref{eq:cauchy-esp}; and the remaining cases follow from
Lemma~\ref{lem:centered-esp}.

For $0\le a\le27/100$,
\[
 \prod_{k=1}^{10}r_k\le\left(\frac{127}{100}\right)^{10}<11.
\]
On $[27/100,37/100]$, Lemma~\ref{lem:sigma} with $M=137/100$ has
$\nu=10$ and gives
\[
 S_0=9\left(\frac{100}{137}\right)^2
 +\left(\frac{(137/100)^9}{11}\right)^2.
\]
The quadratic $S_0a^2-10a+10-S_0$ has positive Bernstein coefficients on
this interval, so \eqref{eq:localization-general} gives $u>1$.

For $i=37,\ldots,89$, define
$\alpha_i,\beta_i,S_i,\lambda_i,Y_i$ exactly as in the degree-ten proof,
with $N=10$, $m=m_{11}$, and $C=11$.  Divide the vertical interval
$[0,Y_i]$ into ten equal parts.  On each of the resulting $530$ rectangles,
use Proposition~\ref{prop:split-certificate} with one value from the fixed
menu
\[
 \mathcal T=\left\{
 1,\frac{19}{20},\frac9{10},\frac{17}{20},\frac45,\frac34,
 \frac7{10},\frac{13}{20},\frac35,\frac{11}{20},\frac12
 \right\}.
\]
For every rectangle at least one choice of $\tau\in\mathcal T$ gives a
polynomial with all tensor-product Bernstein coefficients positive.  The
polynomials have bidegree at most $(11,10)$.  The least coefficient among all
$530$ certificates is greater than $1/2800$.

For $9/10\le a\le1$, take the bound $S_b$ from
Lemma~\ref{lem:sigma} with $M=2$ and set
\[
 \lambda_b=\frac{17}{20}.
\]
Localization gives
\[
 1-u\le C_b(1-a),
 \qquad
 C_b=\frac{S_b(1+9/10)-10}{10(9/10)},
\]
so $\eta^2\le2C_b(1-a)$.  Put $1-a=x^2$ and $\eta=xY$.  The exact
inequalities
\[
 \left(\frac8{25}\right)^2\ge\frac1{10},
 \qquad
 \left(\frac{17}{10}\right)^2\ge2C_b
\]
contain every feasible pair in
\[
 0\le x\le\frac8{25},
 \qquad
 0\le Y\le\frac{17}{10}.
\]
Use the unsplit polynomial \eqref{eq:Pfull}; after division by $x^2$, its
exact Bernstein coefficients are positive on this box;

the least exceeds $3/50$.  This proves degree eleven.

\section{Exact verification}

If
\[
 R(U,V)=\sum_{k=0}^r\sum_{\ell=0}^s a_{k\ell}U^kV^\ell,
\]
then its tensor-product Bernstein coefficients of bidegree $(r,s)$ are
\begin{equation}\label{eq:bern-conversion}
 b_{ij}=
 \sum_{k\le i}\sum_{\ell\le j}a_{k\ell}
 \frac{\binom ik}{\binom rk}
 \frac{\binom j\ell}{\binom s\ell}.
\end{equation}
The ancillary verifiers construct every integral in
\eqref{eq:Pfull} and \eqref{eq:Psplit} symbolically, apply the affine map from
the relevant rational rectangle to $[0,1]^2$, use
\eqref{eq:bern-conversion}, and compare rational numbers exactly.  No
floating-point arithmetic or numerical quadrature enters any certificate.

For each degree there are two independent implementations.  The primary
implementation uses Python integers and rational fractions; the second
reconstructs the polynomials independently in SymPy.  The finite verification
is summarized in Table~\ref{tab:summary}.

\begin{table}[ht]
\centering
\small
\caption{Exact certificate summary.}\label{tab:summary}
\begin{tabular}{@{}c c c c c@{}}
\toprule
Degree & Interior region & Rectangles & Interior lower bound & Boundary lower bound\\
\midrule
9  & $[9/20,9/10]$ & $18$  & $1/2100$ & $1/100$\\
10 & $[2/5,9/10]$  & $50$  & $1/3100$ & $7/100$\\
11 & $[37/100,9/10]$ & $530$ & $1/2800$ & $3/50$\\
\bottomrule
\end{tabular}
\end{table}

The certificate programs and a verification log are included with the source
of this paper.

\section{Conclusion}

The reciprocal normalization, localization inequality, and integral identity
reduce Sendov's conjecture in degrees nine, ten, and eleven to exact
low-dimensional polynomial inequalities.  The degree-ten proof is obtained
from the degree-nine framework by sharper centered symmetric-function bounds.
Degree eleven requires only one additional estimate, namely the split
arithmetic--geometric mean bound in
Proposition~\ref{prop:split-certificate}.  Theorem~\ref{thm:main} follows.

\begin{thebibliography}{9}

\bibitem{BrandenWagner}
P.~Br\"and\'en and D.~G.~Wagner,
\newblock A converse to the Grace--Walsh--Szeg\H{o} theorem,
\newblock \emph{Math. Proc. Cambridge Philos. Soc.} \textbf{147} (2009),
447--453.

\bibitem{BrownXiang}
J.~E.~Brown and G.~Xiang,
\newblock Proof of the Sendov conjecture for polynomials of degree at most
 eight,
\newblock \emph{J. Math. Anal. Appl.} \textbf{232} (1999), 272--292.

\end{thebibliography}

\end{document}
